\documentclass[12pt,a4paper]{article}
\usepackage{fontspec}
\usepackage{polyglossia}
\setmainfont{TeX Gyre Termes}
\newfontfamily\cyrillicfont{TeX Gyre Termes}
\setdefaultlanguage{ukrainian}
\setotherlanguage{english}
\usepackage{amsmath,amssymb,geometry,graphicx}
\geometry{margin=25mm}
\setlength{\parindent}{1.25cm}
\setlength{\parskip}{0.25em}
\begin{document}
\title{2024 MODIFICATION OF THE ALGORITHM FOR DEFINING POLYGONAL GEOMETRY OF AN OBJECT FOR POLYPOINT TRANSFORMATIONS}
\maketitle
DOI https://doi.org/10.15407/csc.202DOI https://doi.org/10.15407/csc.20222.0.033.xxx.xxx
\par
UDKUDK 514.18  514.18
\par
SydorenkoSydorenko  ІІu.Vu.V.., PhD (Eng. sc.), assistant professor, National Technical University of Ukraine “Igor Sikorsky Kyiv Polytechnic Institute”, 37, , PhD (Eng. sc.), assistant professor, National Technical University of Ukraine “Igor Sikorsky Kyiv Polytechnic Institute”, 37, ProspProsp..  PeremohyPeremohy, Kyiv, Ukraine, 03056., , Kyiv, Ukraine, 03056., ORCIDORCID: : https://orcid.org/0000-0002-1953-041https://orcid.org/0000-0002-1953-041
\par
E–mail: suliko3@ukr.netE–mail: suliko3@ukr.net
\par
HorodetskyiHorodetskyi M.V M.V.., , PhD PhD student, National Technical University of Ukraine “Igor Sikorsky Kyiv Polytechnic Institute”, 37, student, National Technical University of Ukraine “Igor Sikorsky Kyiv Polytechnic Institute”, 37, Prosp.PeremohyProsp.Peremohy, Kyiv, Ukraine, 03056., , Kyiv, Ukraine, 03056., ORCID:ORCID:  https://orcid.org/https://orcid.org/0000-0003-4673-38940000-0003-4673-3894
\par
EE––mailmail: :  HYPERLINK "mailto:horodetskyimykola@gmail.com" horodetskyimykolahorodetskyimykola@@gmailgmail..comcom
\par
МОДИФІКАЦІЯ АЛГОРИТМУ ЗАДАННЯ ПОЛІГОНАЛЬНОЇМОДИФІКАЦІЯ АЛГОРИТМУ ЗАДАННЯ ПОЛІГОНАЛЬНОЇ ГЕОМЕТРІЇ ГЕОМЕТРІЇ ОБ’ЄКТА ДЛЯ ПОЛІТОЧКОВИХ ПЕРЕТВОРЕНЬ ОБ’ЄКТА ДЛЯ ПОЛІТОЧКОВИХ ПЕРЕТВОРЕНЬ
\par
MODIFICATION OF THE ALGORITHM FOR MODIFICATION OF THE ALGORITHM FOR DEFININGDEFINING POLYGONAL GEOMETRY OF AN OBJECT FOR POLY POLYGONAL GEOMETRY OF AN OBJECT FOR POLYPOINT POINT TRANSFORMATIONSTRANSFORMATIONS
\par
У статті розглянуто модифікацію побудови полігонів для відображення процесів деформації за допомогою політочкових перетворень. У статті розглянуто модифікацію побудови полігонів для відображення процесів деформації за допомогою політочкових перетворень. Об’єктом політочкових перетворень є пряма, яка  за умови зміни базису змінює свої коефіцієнти. Якщо потрібно провести відображення двовимірної фігури, то цю фігуру потрібно представити як набір відрізків, тобто сформувати полігон. В статті наведено спосіб формування полігону для подальшої його обробки методом Об’єктом політочкових перетворень є пряма, яка  за умови зміни базису змінює свої коефіцієнти. Якщо потрібно провести відображення двовимірної фігури, то цю фігуру потрібно представити як набір відрізків, тобто сформувати полігон. В статті наведено спосіб формування полігону для подальшої його обробки методом полікоординатних полікоординатних перетворень. перетворень.
\par
Ключові слова: Ключові слова: полікоординатні відображення, політочкові перетворення, полігонполікоординатні відображення, політочкові перетворення, полігональна геометріяальна геометрія, базис перетворення, об’єкт перетворення. , базис перетворення, об’єкт перетворення.
\par
KeywordsKeywords::  polycoordinate polycoordinate mappingsmappings, polypoint transformations, polygon, polypoint transformations, polygonal geometryal geometry, transformation basis, transformation object., transformation basis, transformation object.
\par
ВступВступ
\par
На практиці інженери часто стикаються з необхідністю перевірки технічного стану інженерних об’єктів на наявність деформацій. Це потрібно робити не тільки для забезпечення надійності обладнання, а і для прогнозування граничного терміну безвідмовної роботи деяких елементів та окремих вузлів конструкцій. У наш час з проблемами деформацій доволі часто стикаються і представники інших професій, наприклад, лікарі. Під час воєнних дій багато людей потребують відтворення На практиці інженери часто стикаються з необхідністю перевірки технічного стану інженерних об’єктів на наявність деформацій. Це потрібно робити не тільки для забезпечення надійності обладнання, а і для прогнозування граничного терміну безвідмовної роботи деяких елементів та окремих вузлів конструкцій. У наш час з проблемами деформацій доволі часто стикаються і представники інших професій, наприклад, лікарі. Під час воєнних дій багато людей потребують відтворення пошкоджених, деформованих органів. Інколи це може врятувати життя пацієнту. Необхідність розв’язувати задачі деформаційного моделювання виникають також при прогнозуванні динаміки розповсюдження різних забруднень, наприклад при застосуванні хімічної зброї, або для розуміння, куди пошириться пожежа при «прильотах», для швидкого реагування на ситуацію.  пошкоджених, деформованих органів. Інколи це може врятувати життя пацієнту. Необхідність розв’язувати задачі деформаційного моделювання виникають також при прогнозуванні динаміки розповсюдження різних забруднень, наприклад при застосуванні хімічної зброї, або для розуміння, куди пошириться пожежа при «прильотах», для швидкого реагування на ситуацію.
\par
Одним з методів розв’язання задачі деформаційного моделювання є політочкові перетворення. Апарат цих перетворень допомагає швидко отримувати математичну модель перетворень в залежності від наслідків деформації. Кожен об’єкт представляється набором точок, які  користувач задає в процесі роботи алгоритму. Ці точки є перетином прямих. На виході необхідно отримати знову таки набір точок, що будуть формувати перетворений об’єкт. Але існує проблема в отриманні точок перетинуОдним з методів розв’язання задачі деформаційного моделювання є політочкові перетворення. Апарат цих перетворень допомагає швидко отримувати математичну модель перетворень в залежності від наслідків деформації. Кожен об’єкт представляється набором точок, які  користувач задає в процесі роботи алгоритму. Ці точки є перетином прямих. На виході необхідно отримати знову таки набір точок, що будуть формувати перетворений об’єкт. Але існує проблема в отриманні точок перетину через те, що геометрія об’єкта задається багатьма прямими через те, що геометрія об’єкта задається багатьма прямими, тому і точок перетину буде багато. І потрібно визначити, які саме точки перетину утворюють змінений об’єкт. Модифікація способу задання об’єкта дозволить швидко отримувати ці точки., тому і точок перетину буде багато. І потрібно визначити, які саме точки перетину утворюють змінений об’єкт. Модифікація способу задання об’єкта дозволить швидко отримувати ці точки.
\par
Аналіз останніх дослідженьАналіз останніх досліджень. . У роботіУ роботі  [1] наведено постановку задачі та методи реалізації політочкових перетворень. Використання цих перетворень набуло розвитку у представлені об’єктів різними способами, що описано у роботі[1] наведено постановку задачі та методи реалізації політочкових перетворень. Використання цих перетворень набуло розвитку у представлені об’єктів різними способами, що описано у роботі  [2]. [2].
\par
При роботі з політочковими перетвореннями виникає необхідність отримати форму об’єкта після перетворення. В результаті роботи алгоритму на виході є набір точок, який потрібно поєднати у неперервну криву. Це можна зробити у різний спосіб, наприклад, застосовуючи інтерполяційну криву Гауса, як описано у роботі[3]. При роботі з політочковими перетвореннями виникає необхідність отримати форму об’єкта після перетворення. В результаті роботи алгоритму на виході є набір точок, який потрібно поєднати у неперервну криву. Це можна зробити у різний спосіб, наприклад, застосовуючи інтерполяційну криву Гауса, як описано у роботі[3].
\par
Нові варіанти розв’язання задачі моделювання деформованих об’єктів дозволило підвищити точність розрахунків отриманих точок. Це було досягнуто за рахунок введення нового функціоналу, що описано у роботі [4].Нові варіанти розв’язання задачі моделювання деформованих об’єктів дозволило підвищити точність розрахунків отриманих точок. Це було досягнуто за рахунок введення нового функціоналу, що описано у роботі [4].
\par
Розробки у напрямку подальшого розвитку полікоординатних відображень дозволили розв’язати низку практичних задач, таких як підрахунок площі замкненого контуруРозробки у напрямку подальшого розвитку полікоординатних відображень дозволили розв’язати низку практичних задач, таких як підрахунок площі замкненого контуру  [4], візуалізація масових сцен в комп’ютерних іграх[4], візуалізація масових сцен в комп’ютерних іграх  [4], моделювання конструктивних елементів вентиляційних систем[4], моделювання конструктивних елементів вентиляційних систем  [4], моделювання та візуалізація деформованих органів у пластичній хірургії[4], моделювання та візуалізація деформованих органів у пластичній хірургії  [4] та інших напрямках. [4] та інших напрямках.
\par
Ефективність процесу відображення деформації за допомогою політочкових перетворень залежить від способу представлення об’єкта перетворення. У даній роботі розглядаєтьсЕфективність процесу відображення деформації за допомогою політочкових перетворень залежить від способу представлення об’єкта перетворення. У даній роботі розглядається алгоритм представлення початкового об’єкта у вигляді набору полігонів, заданих певним чином, що спрощує відображення цього ж об’єкта після застосування до нього політочкових перетворень.я алгоритм представлення початкового об’єкта у вигляді набору полігонів, заданих певним чином, що спрощує відображення цього ж об’єкта після застосування до нього політочкових перетворень.
\par
Метою дослідженняМетою дослідження є є модифікація алгоритму політочкових перетворень за рахунок  модифікація алгоритму політочкових перетворень за рахунок   розробки способу розробки способу зазаввдання дання точокточок об’єкта і об’єкта і розроб розробки гнучкогоки гнучкого підх підхооддуу  на основі на основі застосування застосування полігонівполігонів певного вигляду певного вигляду..
\par
Виклад основного матеріалуВиклад основного матеріалу..  Політочкові перетворення є одним з різновидів полікоординатних відображень, які використовуються для розв’язання задач відображення деформаційних змін геометричних об’єктів. Політочкові перетворення є одним з різновидів полікоординатних відображень, які використовуються для розв’язання задач відображення деформаційних змін геометричних об’єктів.
\par
Об’єктом деформації при Об’єктом деформації при пполіточкових перетвореннях у двовимірному просторі (Rоліточкових перетвореннях у двовимірному просторі (R22))  є прямає пряма  такого виглядутакого вигляду::
\par
Нехай задана пряма Нехай задана пряма LL11, з коефіцієнтами , з коефіцієнтами AA11, B, B11, C, C11. . Розв’язання задачі Розв’язання задачі політочкових політочкових перетворень зводиться до розв’язання перетворень зводиться до розв’язання системи лінійних алгебричних рівнянь,системи лінійних алгебричних рівнянь, яку можна записати так: яку можна записати так:
\par
де де XXіі,Y,Yіі — координата  — координата іі-того базисного вузла,-того базисного вузла,  ββii — відстань від прямої  — відстань від прямої LL1 1 до базиса до базиса XXіі,Y,Yіі  ..
\par
Розв’язкамиРозв’язками системи  системи є три коефіцієнти є три коефіцієнти A, B, CA, B, C  , , підставивши які у загальне рівняння прямоїпідставивши які у загальне рівняння прямої,,  отримаємо нову перетворену пряму отримаємо нову перетворену пряму LL.. На рисунку 1 представлено політочкові перетворення прямої.  На рисунку 1 представлено політочкові перетворення прямої.
\par
РисРис. . 11. . Політочкові перетіорення прямої Політочкові перетіорення прямої LL1 1 з використанням трьохз використанням трьох базисних точок базисних точок
\par
Оскільки за допомогою політочкових перетворень можна перенести з одного базису в другий одну пряму, то можна і декілька. Для прикладу побудуємо сітку з прямих, яка являє собою координатну сітку. Це буде множина ортогональних прямих з рівним кроком. Керуючи точками базису отримаємо змінену координатну сітку, як показано на рисунку 2. Оскільки за допомогою політочкових перетворень можна перенести з одного базису в другий одну пряму, то можна і декілька. Для прикладу побудуємо сітку з прямих, яка являє собою координатну сітку. Це буде множина ортогональних прямих з рівним кроком. Керуючи точками базису отримаємо змінену координатну сітку, як показано на рисунку 2.
\par
Рис. 2. Перетворення координатної сіткиРис. 2. Перетворення координатної сітки
\par
Перетворення довільних двовимірних геометричних об’єктів можна проводити різними способами. Можна використати ту саму координатну сітку. Відобразити на ній об’єкт, розбити його на точки, а потім шукати на зміненій сітці відповідні точки. Цей процес ускладняється тим, що на виході при роботі апарату політочкових перетворень результатом є коефіцієнти прямих. Тобто, деякі набори векторів з трьох чисел. І зрозуміти, де саме буде точка перетину двох конкретних прямихПеретворення довільних двовимірних геометричних об’єктів можна проводити різними способами. Можна використати ту саму координатну сітку. Відобразити на ній об’єкт, розбити його на точки, а потім шукати на зміненій сітці відповідні точки. Цей процес ускладняється тим, що на виході при роботі апарату політочкових перетворень результатом є коефіцієнти прямих. Тобто, деякі набори векторів з трьох чисел. І зрозуміти, де саме буде точка перетину двох конкретних прямих,, доволі складно.  доволі складно.
\par
Є інший підхід, який часто застосовується при роботі з політочковими перетвореннями. Кожна точка об’єкта подається у вигляді перетину двох конкретних ортогональних прямих, і тоді нема потреби відслідковувати всі перетини всіх заданих прямих. Є інший підхід, який часто застосовується при роботі з політочковими перетвореннями. Кожна точка об’єкта подається у вигляді перетину двох конкретних ортогональних прямих, і тоді нема потреби відслідковувати всі перетини всіх заданих прямих. Переносяться дві Переносяться дві прямпряміі у новий базис, знаходиться точка перетину. І так  у новий базис, знаходиться точка перетину. І так по кожній точціпо кожній точці з кожною точкою з кожною точкою. . Але у цьому підході Але у цьому підході єє той той недолік недолік: : , що, що оскільки точка представляється перетином двох прямих, то кількість обчислень зростає вдвічі, що в свою чергу  оскільки точка представляється перетином двох прямих, то кількість обчислень зростає вдвічі, що в свою чергу призводить до збільшення часу обчислень призводить до збільшення часу обчислень а такожа також використання зайвої пам’яті комп’ютера.  використання зайвої пам’яті комп’ютера.
\par
УУ даному дослідженні  даному дослідженні розглядається спосіброзглядається спосіб  заданнязадання  об’єкта за допомогою об’єкта за допомогою відрізківвідрізків, які з’єднують сусідні точки, утворюючи полігон. , які з’єднують сусідні точки, утворюючи полігон.
\par
При реалізації такого підходу потрібно додержуватись певних правил. При реалізації такого підходу потрібно додержуватись певних правил.
\par
По-перше, По-перше, відрізки повинні бути поєднанівідрізки повинні бути поєднані, тобто кінець одного відрізка по, тобто кінець одного відрізка повинен бути початком наступного винен бути початком наступного  (рис. 3). (рис. 3).
\par
Рис. 3. Об’єкт, заданий непоєднаними та поєднаними відрізкамиРис. 3. Об’єкт, заданий непоєднаними та поєднаними відрізками
\par
По-друге, даний підхід вимагаєПо-друге, даний підхід вимагає, щоб об’єкт був замкнений, щоб об’єкт був замкнений. Тобто геометрія об’єкта має задаватись шляхом поєднання відрізків один за одни. Тобто геометрія об’єкта має задаватись шляхом поєднання відрізків один за однимм, утворюючи ланцюжок з відрізків (шлях), кінцем якого, утворюючи ланцюжок з відрізків (шлях), кінцем якого завжди повинен бути його почат завжди повинен бути його початок.ок.
\par
Таким чином, алгоритм перетворення геометричного об’єкта буде виглядати наступним чином:Таким чином, алгоритм перетворення геометричного об’єкта буде виглядати наступним чином:
\par
ЗадаємоЗадаємо  ВводимоВводимо початковий базис перетворень (певну кількість точок).  початковий базис перетворень (певну кількість точок).
\par
ЗадаємоЗадаємо  ВводимоВводимо  послідовно точки відрізків, на які розбито послідовно точки відрізків, на які розбито перетворюємий перетворюємий об’єкт. Точки кінця попереднього відрізку і початку наступного співпадаютьоб’єкт. Точки кінця попереднього відрізку і початку наступного співпадають  іі  тожтож їх не потрібно дублювати.  їх не потрібно дублювати. Вони враховуються автоматично.Вони враховуються автоматично.
\par
Точки потрапляютьТочки потрапляють у чергу типу  у чергу типу FIFOFIFO  LIFOLIFO, властивістю якої є обробка кожного , властивістю якої є обробка кожного даногоданого  елементуелементу у порядку входження у порядку входження..
\par
По По кожнкожній зій з двох послідов двох послідовнноо заданих точках будуються прямі, які є об’єктами перетворень. заданих точках будуються прямі, які є об’єктами перетворень.
\par
Проводиться зміна початкового базису шляхом переміщення базисних точок.Проводиться зміна початкового базису шляхом переміщення базисних точок.
\par
Проходить розрахунок за методом політочкових перетворень нових положень прямих.Проходить розрахунок за методом політочкових перетворень нових положень прямих.  Проводяться розрахунки методом політочкових перетворень нових положень прямих.Проводяться розрахунки методом політочкових перетворень нових положень прямих.
\par
На виході отримуємо точки перетину цих прямих, що і будуть утворювати форму перетвореного об’єкта.На виході отримуємо точки перетину цих прямих, що і будуть утворювати форму перетвореного об’єкта.
\par
Проводимо візуалізацію Проводимо візуалізацію знайденогознайденого контуру контуру  використвикориставшиавши той самий порядок точок що і в П3 той самий порядок точок що і в П3..  видаляючи ті фрагментвидаляючи ті фрагментии прямих,  прямих, щощо не є частиною об’єкта не є частиною об’єкта  . .
\par
Розглянемо це на прикладі.Розглянемо це на прикладі.  Нехай потрібно задати об’єкт у двовимірному просторі за допомогою відрізків (далі ребНехай потрібно задати об’єкт у двовимірному просторі за допомогою відрізків (далі реберер) для подальшого застосування ) для подальшого застосування методу пметоду політочкових перетворень. оліточкових перетворень.
\par
Для цього кожне ребро задаДля цього кожне ребро задамо мо послідовно у напрямку слідування контуру об’єкта використовуючи дві точки послідовно у напрямку слідування контуру об’єкта використовуючи дві точки (x(xii; y; yii), (x), (xjj; y; yjj)) для кожного. Пронумирувавши кожен кінець кожного ребра отримаємо закритий полігон  для кожного. Пронумирувавши кожен кінець кожного ребра отримаємо закритий полігон ««ABCDEFGHAABCDEFGHA»» (далі AВ-HA) (рис.  (далі AВ-HA) (рис. 44). Відмітимо, що останнє ребро НА закінчується на початку ребра АВ.). Відмітимо, що останнє ребро НА закінчується на початку ребра АВ.
\par
Рис. 4Рис. 4. . Приклад замкненого Приклад замкненого пполігонолігонуу
\par
Перетворимо полігон в прямі Перетворимо полігон в прямі загального виглядузагального вигляду, застосуємо п, застосуємо політочкові перетворенняоліточкові перетворення,, та знайдемо точки перетину сусідніх прямих стандартним та знайдемо точки перетину сусідніх прямих стандартним способом.  способом.
\par
Отримаємо Отримаємо нове положення нове положення прямихпрямих у перетвореному базисі у перетвореному базисі (рис.5) (рис.5). .
\par
РисРис. . 55. . ПоложенняПоложення  прямих прямих у перетвореному базисіу перетвореному базисі
\par
Після цього на екрані Після цього на екрані комп’ютеракомп’ютера видаляємо ті частини прямих, які не утворюють полігон (знаходяться за межами полігону) видаляємо ті частини прямих, які не утворюють полігон (знаходяться за межами полігону)..  Використавши точки перетинуВикориставши точки перетину  між сусіднми прямимиміж сусіднми прямими  прямими із списку прямипрямими із списку прямихх утвор утворимоимо новий полігон  новий полігон «A«A11BB11CC11DD11EE11FF11GG11HH11AA11»»  (далі (далі AA11BB11- H- H11AA11) , який ) , який єє п праобразом полігону AВ-HA (рис. 6раобразом полігону AВ-HA (рис. 6).).
\par
РисРис. . 66. . Полігон після політочкових перетвореньПолігон після політочкових перетворень
\par
Таким чиномТаким чином, було отримано перетворений об’єкт після деформації.  П, було отримано перетворений об’єкт після деформації.  Перетворений полігон еретворений полігон «A«A11BB11CC11DD11EE11FF11GG11HH11AA11»», складений із тією самою кількістю ребер що і початковий полігон , складений із тією самою кількістю ребер що і початковий полігон «ABCDEFGHA», який «ABCDEFGHA», який побудовано за законом зміни точок базису.побудовано за законом зміни точок базису.
\par
Такий же підхід можна використати і для зважених політочкових перетворень. Для цього потрібно застосувати інший функціонал, а саме функціонал такого вигляду:Такий же підхід можна використати і для зважених політочкових перетворень. Для цього потрібно застосувати інший функціонал, а саме функціонал такого вигляду:
\par
,,
\par
де де mm((ββii)) –  – функція функція від відстаней до прямої образу. Наприклад, нехай від відстаней до прямої образу. Наприклад, нехай mm буде обернено-пропорційно залежати від  буде обернено-пропорційно залежати від ββ у парній степені. У цьому випадку точки об’єкта будуть тим більше притягуватись до точок базису, чим ближче вони до цих точок знаходяться. Причому, чим більша степінь, тим сильніший зв'язок.  у парній степені. У цьому випадку точки об’єкта будуть тим більше притягуватись до точок базису, чим ближче вони до цих точок знаходяться. Причому, чим більша степінь, тим сильніший зв'язок.
\par
Також на екрані можна розмістити декілька об’єктів, і застосувати описаний підхід до задання Також на екрані можна розмістити декілька об’єктів, і застосувати описаний підхід до задання кожного кожного об’єкта. об’єкта. Наприклад, задаємо 4 об’єкта на одній сценіНаприклад, задаємо 4 об’єкта на одній сцені (рис.7 (рис.7). ).
\par
РисРис. 7. 7. Об’єкт з чотирьох полігонів після . Об’єкт з чотирьох полігонів після пполіточкових перетворень оліточкових перетворень
\par
Як видно з рисункуЯк видно з рисунку 7 7, всі чотири об’єкти відобразились коректно. Причому, швидкість відображення об’єктів , всі чотири об’єкти відобразились коректно. Причому, швидкість відображення об’єктів співспіврозрозмірнамірна з їх кількістю і обчислюється долями мікросекунд, на відміну від відображення результатів без застосування полігону.  з їх кількістю і обчислюється долями мікросекунд, на відміну від відображення результатів без застосування полігону.
\par
ВисновкиВисновки..  РозглянутоРозглянуто проблеми побудови двовимірних об’єктів для подальш проблеми побудови двовимірних об’єктів для подальшого застосування до них методу пого застосування до них методу політочкових перетворень. Запропоновано дієвий метод полігонального задання об’єкта у двовимірному просторі для збереження початкової кількості відрізків у перетвореному праобразі об’єкта. Продемонстровано використання представленого алгоритму задання об’єкту для зображення оліточкових перетворень. Запропоновано дієвий метод полігонального задання об’єкта у двовимірному просторі для збереження початкової кількості відрізків у перетвореному праобразі об’єкта. Продемонстровано використання представленого алгоритму задання об’єкту для зображення декількох об’єктівдекількох об’єктів на одній сцені. на одній сцені.
\par
ЛІТЕРАТУРА ЛІТЕРАТУРА
\par
Сидоренко Ю.В. Система моделювання геометричних об’єктів за допомогою політочкових перетворень. Прикладна геометрія та інженерна графіка. - К.:КДТУБА, 2016. - Вип.92- с.118-125.Сидоренко Ю.В. Система моделювання геометричних об’єктів за допомогою політочкових перетворень. Прикладна геометрія та інженерна графіка. - К.:КДТУБА, 2016. - Вип.92- с.118-125.
\par
Сидоренко Ю.В., Бадаєв Ю.І. Геометричне моделювання складних об’єктів на основі політочкових відображень відрізків прямих. Сидоренко Ю.В., Бадаєв Ю.І. Геометричне моделювання складних об’єктів на основі політочкових відображень відрізків прямих. СучасніСучасні  проблемипроблеми  моделюваннямоделювання. 2019. -Вип.16.-С.17-24.  . 2019. -Вип.16.-С.17-24.
\par
СидоренкоСидоренко Ю.В.,  Ю.В., ШалденкоШалденко О.В.  О.В. ВізуалізаціяВізуалізація  об’єктівоб’єктів  політочковихполіточкових  перетвореньперетворень  заза  допомогоюдопомогою  інтерполяційноїінтерполяційної  функціїфункції Гауса. Сучасні проблеми моделювання. 2020.- Вип.17.-С.108-114.  Гауса. Сучасні проблеми моделювання. 2020.- Вип.17.-С.108-114.
\par
Сидоренко Ю.В., Залевська О.В. Підвищення точності алгоритму політочкових перетворень.  Сидоренко Ю.В., Залевська О.В. Підвищення точності алгоритму політочкових перетворень.  ПрикладнаПрикладна  геометріягеометрія  тата  інженернаінженерна графіка.  2020, вип. графіка.  2020, вип.97  —97  —  С.129-135.      С.129-135.
\par
Сидоренко Ю.В., Залевська О.В., Шалденко О. В. Обчислення площі перетвореного об’єкта при політочкових перетвореннях. Сидоренко Ю.В., Залевська О.В., Шалденко О. В. Обчислення площі перетвореного об’єкта при політочкових перетвореннях. ПрикладнаПрикладна  геометріягеометрія  тата  інженернаінженерна графіка. 2022. Випуск 102. Київ. Ст.65-75  графіка. 2022. Випуск 102. Київ. Ст.65-75
\par
Сидоренко, Ю., Шалденко, О., \& Онисько, А. (2023). Система політочкових перетворень групи об’єктів в одному базисі. Сидоренко, Ю., Шалденко, О., \& Онисько, А. (2023). Система політочкових перетворень групи об’єктів в одному базисі. СучасніСучасні  проблемипроблеми  моделюваннямоделювання, (25), 202-209  , (25), 202-209
\par
Sidorenko Yu, Sidorenko Yu, KryvdaKryvda O.,  O., LeshchynskaLeshchynska I. System of modeling of structural elements of ventilation systems by  I. System of modeling of structural elements of ventilation systems by polycoоrdinatepolycoоrdinate transformations. Strength of Materials and Theory of Structures. 2020. No. 104. P. 221-228.  transformations. Strength of Materials and Theory of Structures. 2020. No. 104. P. 221-228.
\par
СидоренкоСидоренко Ю.В.,  Ю.В., ШалденкоШалденко О.В.  О.В. ВаговіВагові  політочковіполіточкові  перетворенняперетворення  припри  моделюваннімоделюванні  можливихможливих  результатіврезультатів  пластичноїпластичної хірургії. Сучасні проблеми моделювання. 2019. – Вип.15. – С. 151-161. хірургії. Сучасні проблеми моделювання. 2019. – Вип.15. – С. 151-161.
\par
REFERENCESREFERENCES
\par
1.1.	SidorenkoSidorenko  YuYu..VV. . SistemaSistema  modelyuvannyamodelyuvannya  geometrichnikhgeometrichnikh  obob''yektiyekti``vv  zaza  dopomogoyudopomogoyu  polipoli``tochkovikhtochkovikh  peretvorenperetvoren`. `. PrikladnaPrikladna  geometri`yageometri`ya ta  ta i`nzhenernai`nzhenerna  grafi`kagrafi`ka. - . - K.:KDTUBAK.:KDTUBA, 2016. - Vip.92- s.118-125., 2016. - Vip.92- s.118-125.
\par
2.2.	Sidorenko Sidorenko Yu.VYu.V., ., BadayevBadayev  Yu.IYu.I`. `. GeometrichneGeometrichne  modelyuvannyamodelyuvannya  skladnikhskladnikh  ob'yekti`vob'yekti`v  nana  osnoviosnovi` ` poli`tochkovikhpoli`tochkovikh  vi`dobrazhenvi`dobrazhen` ` vi`dri`zki`vvi`dri`zki`v  pryamikhpryamikh. . SuchasniSuchasni` ` problemiproblemi  modelyuvannyamodelyuvannya. 2019. -Vip.16.-S.17-24.  . 2019. -Vip.16.-S.17-24.
\par
3.3.	Sidorenko Sidorenko Yu.VYu.V., ., ShaldenkoShaldenko O.V.  O.V. Vi`zuali`zaczi`yaVi`zuali`zaczi`ya  ob'yekti`vob'yekti`v  poli`tochkovikhpoli`tochkovikh  peretvorenperetvoren` za ` za dopomogoyudopomogoyu  i`nterpolyaczi`jnoyii`nterpolyaczi`jnoyi  funkczi`yifunkczi`yi  GausaGausa. . SuchasniSuchasni` ` problemiproblemi  modelyuvannyamodelyuvannya. 2020.- Vip.17.-S.108-114. . 2020.- Vip.17.-S.108-114.
\par
4.4.	Sidorenko Sidorenko Yu.VYu.V., ., Zalevs`kaZalevs`ka O.V.  O.V. Pi`dvishhennyaPi`dvishhennya  tochnostitochnosti` ` algoritmualgoritmu  poli`tochkovikhpoli`tochkovikh  peretvorenperetvoren`.  `.  PrikladnaPrikladna  geometri`yageometri`ya ta  ta i`nzhenernai`nzhenerna  grafi`kagrafi`ka.  2020, vip..  2020, vip.97  —97  —  S.129-135.      S.129-135.
\par
5.5.	Sidorenko Sidorenko Yu.VYu.V., ., Zalevs`kaZalevs`ka O.V.,  O.V., ShaldenkoShaldenko O. V.  O. V. ObchislennyaObchislennya  ploshhiploshhi` ` peretvorenogoperetvorenogo  ob'yektaob'yekta  pripri  poli`tochkovikhpoli`tochkovikh  peretvorennyakhperetvorennyakh. . PrikladnaPrikladna  geometri`yageometri`ya ta  ta i`nzhenernai`nzhenerna  grafi`kagrafi`ka. 2022. . 2022. VipuskVipusk 102.  102. KiyivKiyiv. St.65-75 . St.65-75
\par
6.6.	Sidorenko, Yu., Sidorenko, Yu., ShaldenkoShaldenko, O., \& , O., \& Onis`koOnis`ko, A. (2023). Sistema , A. (2023). Sistema poli`tochkovikhpoli`tochkovikh  peretvorenperetvoren` ` grupigrupi  ob'yekti`vob'yekti`v v  v odnomuodnomu  bazisibazisi`. `. SuchasniSuchasni` ` problemiproblemi  modelyuvannyamodelyuvannya, (25), 202-209  , (25), 202-209
\par
7.7.	Sidorenko Yu, Sidorenko Yu, KryvdaKryvda O.,  O., LeshchynskaLeshchynska I. System of modeling of structural elements of ventilation systems by polycoordinate transformations. Strength of Materials and Theory of Structures. 2020. No. 104. P. 221-228.  I. System of modeling of structural elements of ventilation systems by polycoordinate transformations. Strength of Materials and Theory of Structures. 2020. No. 104. P. 221-228.
\par
8. 8. Sidorenko Sidorenko Yu.VYu.V., ., ShaldenkoShaldenko O.V.  O.V. VagoviVagovi` ` poli`tochkovipoli`tochkovi` ` peretvorennyaperetvorennya  pripri  modelyuvannimodelyuvanni` ` mozhlivikhmozhlivikh  rezul`tati`vrezul`tati`v  plastichnoyiplastichnoyi  khi`rurgi`yikhi`rurgi`yi. . SuchasniSuchasni` ` problemiproblemi  modelyuvannyamodelyuvannya. 2019. – Vip.15. – S. 151-161.. 2019. – Vip.15. – S. 151-161.
\par
Ю.В. СИДОРЕНКОЮ.В. СИДОРЕНКО,,  канд.канд.  техн.техн.  наук,наук, доцент,  доцент, Національний технічний університет Національний технічний університет 
України «Київський політехнічний інститут імені Ігоря Сікорського», України «Київський політехнічний інститут імені Ігоря Сікорського», 
03056, м. Київ, просп. Перемоги, 37, Україна, 03056, м. Київ, просп. Перемоги, 37, Україна, ORCIDORCID: : httpshttps://://orcidorcid..orgorg// HYPERLINK "https://orcid.org/0000-0002-1953-041" 0000-0002-1953-0410000-0002-1953-041      
sulikosuliko3@3@ukrukr..netnet
\par
М.В. ГОРОДЕЦЬКИЙМ.В. ГОРОДЕЦЬКИЙ,,  аспірантаспірант, , Національний технічний університет України Національний технічний університет України 
«Київський політехнічний інститут імені Ігоря Сікорського», «Київський політехнічний інститут імені Ігоря Сікорського», 
03056, м. Київ, просп. Перемоги, 37, Україна, 03056, м. Київ, просп. Перемоги, 37, Україна, ORCIDORCID: : httpshttps://://orcidorcid..orgorg//0000-0003-4673-38940000-0003-4673-3894
horodetskyimykolahorodetskyimykola@@gmailgmail..comcom
\par
МОДИФІКАЦІЯ АЛГОРИТМУ ЗАДАННЯ ПОЛІГОНАЛЬНОЇ ГЕОМЕТРІЇ ОБ’ЄКТА ДЛЯ ПОЛІТОЧКОВИХ ПЕРЕТВОРЕНЬМОДИФІКАЦІЯ АЛГОРИТМУ ЗАДАННЯ ПОЛІГОНАЛЬНОЇ ГЕОМЕТРІЇ ОБ’ЄКТА ДЛЯ ПОЛІТОЧКОВИХ ПЕРЕТВОРЕНЬ
\par
Вступ. Вступ. У наш час У наш час інженериінженери вимушені  вимушені регулярно проводрегулярно проводити ити перевірку технічного стану об'єктів для виявлення деформацій, що є важливим для забезпечення надійноперевірку технічного стану об'єктів для виявлення деформацій, що є важливим для забезпечення надійної, безвідмовної ї, безвідмовної та та прогнозованої прогнозованої роботи елементів конструкцій. Проблеми деформацій також актуальні у медицині, де відтворення пошкоджених органів може бути життєво важливим, особливо під час воєнних конфліктів. У розв'язанні задач деформаційного моделювання застосовують політочкові перетворення, які дозволяють швидко отримувати математичну модель роботи елементів конструкцій. Проблеми деформацій також актуальні у медицині, де відтворення пошкоджених органів може бути життєво важливим, особливо під час воєнних конфліктів. У розв'язанні задач деформаційного моделювання застосовують політочкові перетворення, які дозволяють швидко отримувати математичну модель перетвореньперетворень у залежності від наслідків деформацій, використовуючи  у залежності від наслідків деформацій, використовуючи поліполіточкове представлення об'єкт.точкове представлення об'єкт.
\par
МетоюМетою дослідження є розібрати існуючі проблеми задання гометрії об’єкта для політочкових перетворень і розробити практичний та гнучкий підхід для їх вирішення на базі полігональної геометрії. дослідження є розібрати існуючі проблеми задання гометрії об’єкта для політочкових перетворень і розробити практичний та гнучкий підхід для їх вирішення на базі полігональної геометрії.
\par
Методи.Методи. Для виконання проведеного дослідження застосовано  Для виконання проведеного дослідження застосовано метод задання геометрії об’єкта через шлях або поліглнметод задання геометрії об’єкта через шлях або поліглн, а також метод п, а також метод політочкових перетворень для отримання видозміненого праобразу вхідної геометрії об’єкта.оліточкових перетворень для отримання видозміненого праобразу вхідної геометрії об’єкта.
\par
РезультатиРезультати..  У статті досліджуютьсяУ статті досліджуються існуючі обмеження щодо зажання геометрії об’єкта на двовимірному просторі та наводиться модифікований алгоритм задання об’єкту для проведення політочкових перетворень на базі полігонів існуючі обмеження щодо зажання геометрії об’єкта на двовимірному просторі та наводиться модифікований алгоритм задання об’єкту для проведення політочкових перетворень на базі полігонів,, а також представлення об’єкта після застосування вказаного методу. а також представлення об’єкта після застосування вказаного методу.
\par
Висновки. Висновки. За результатами дослідження виявилось, що За результатами дослідження виявилось, що модифікований алгоритм задання геометрії об’єкта модифікований алгоритм задання геометрії об’єкта на базі полігонів на базі полігонів є зручнимє зручним, не потребує додаткових обчислень і за , не потребує додаткових обчислень і за допомогою цього підходу можливо задати геометрію для об’єкта будь-якої складності, придатну для застосування мотоду політочкових перетворень.допомогою цього підходу можливо задати геометрію для об’єкта будь-якої складності, придатну для застосування мотоду політочкових перетворень.
\par
Ключові словаКлючові слова: : полікоординатні відображення, політочкові перетворення, полігон, базис перетворення, об’єкт перетворення.полікоординатні відображення, політочкові перетворення, полігон, базис перетворення, об’єкт перетворення.
\par
\end{document}
